\chapter[Modelagem do banco de dados]{Modelagem do banco de dados} \label{capitulo7}

A modelagem do banco de dados foi desenvolvida com o objetivo de estruturar as informações utilizadas no processo de avaliação psicopedagógica, mantendo a organização dos registros e a relação entre aprendentes, profissionais, instrumentos avaliativos, aplicações e relatórios. O banco de dados foi implementado em MySQL e organizado de forma relacional, com utilização de identificadores únicos e chaves estrangeiras para estabelecer as relações entre as entidades.

O modelo contempla inicialmente a entidade \textit{users}, responsável pelo cadastro dos usuários do sistema. Cada usuário possui identificador único, nome, e-mail, senha armazenada por meio de hash, perfil de acesso e status. Foram definidos os perfis \textit{ADMIN}, \textit{PSICOPEDAGOGO} e \textit{VIEWER}, permitindo diferenciar os níveis de acesso às funcionalidades do sistema. A entidade também possui campos para controle de criação, atualização e exclusão lógica dos registros. Associadas aos usuários estão as entidades \textit{refresh\_tokens} e \textit{password\_resets}, utilizadas, respectivamente, para o gerenciamento da renovação das sessões de autenticação e para o processo de recuperação de senha.

A entidade \textit{aprendentes} concentra as informações cadastrais dos indivíduos acompanhados pelo profissional. São armazenados dados como nome, data de nascimento, sexo, informações do responsável, contatos, e-mail e informações relacionadas ao contexto escolar, incluindo escola, série e turno. Também foi disponibilizado um campo para observações. A entidade possui ainda os campos de controle de criação, atualização e exclusão lógica. Para possibilitar o registro longitudinal de informações, a entidade \textit{historico\_aprendentes} mantém registros associados a cada aprendente, contendo uma descrição e uma data, estabelecendo uma relação de um para muitos entre aprendentes e seus históricos.

A representação dos instrumentos utilizados no processo avaliativo foi realizada por meio da entidade \textit{avaliacoes}. Cada avaliação possui nome, descrição, categoria, versão e status, podendo assumir os estados \textit{DRAFT}, \textit{ACTIVE} ou \textit{ARCHIVED}. Um aspecto relevante dessa modelagem é a utilização do campo \textit{parentId}, que estabelece uma autorrelação entre avaliações e permite representar diferentes versões de um mesmo instrumento. Dessa forma, alterações realizadas em um instrumento podem ser mantidas como novas versões, preservando a relação com sua versão anterior.

Os itens que compõem cada avaliação são armazenados na entidade \textit{questions}. Cada questão está vinculada a uma avaliação específica e possui texto, tipo, obrigatoriedade e posição de apresentação. O modelo contempla diferentes tipos de resposta, incluindo texto curto, texto longo, valor numérico, data, seleção única, seleção múltipla, escala e respostas de sim ou não. Para questões de seleção, são armazenadas as opções correspondentes, enquanto questões de escala podem possuir valores mínimo e máximo e respectivos rótulos. Essa estrutura permite que os instrumentos sejam configurados de maneira flexível, sem que cada protocolo precise possuir uma estrutura de banco de dados específica.

A aplicação de um instrumento a um aprendente é representada pela entidade \textit{aplicacoes}. Essa entidade estabelece a relação entre o aprendente, a avaliação utilizada e o usuário responsável pela aplicação. Além dessas associações, são registrados o estado da aplicação, a data de aplicação, observações e a data de finalização. Os estados definidos são \textit{PENDING}, \textit{IN\_PROGRESS}, \textit{DRAFT} e \textit{COMPLETED}, permitindo acompanhar diferentes momentos do preenchimento de uma avaliação. Assim, uma mesma avaliação pode ser aplicada a diferentes aprendentes, enquanto um aprendente pode possuir múltiplas aplicações ao longo de seu acompanhamento.

As respostas fornecidas durante uma aplicação são armazenadas na entidade \textit{respostas}. Cada registro relaciona uma aplicação a uma questão específica e possui um campo \textit{valor} do tipo JSON, permitindo armazenar diferentes formatos de resposta de acordo com o tipo da questão. Foi definida ainda uma restrição de unicidade composta entre \textit{aplicacaoId} e \textit{questionId}, impedindo que uma mesma questão possua mais de uma resposta dentro da mesma aplicação.

Por fim, a entidade \textit{relatorios} representa a devolutiva associada a uma aplicação concluída. Cada relatório está vinculado a uma aplicação específica e possui um autor, além dos campos destinados à conclusão e às observações do profissional. A relação entre \textit{aplicacoes} e \textit{relatorios} é de zero ou um, de modo que uma aplicação pode não possuir relatório inicialmente, mas pode receber uma devolutiva posteriormente. Essa estrutura mantém a elaboração das conclusões vinculada ao registro da avaliação, sem transferir ao sistema a responsabilidade pela interpretação automática dos resultados.

De forma geral, os principais relacionamentos do modelo são apresentados no Diagrama Entidade-Relacionamento. Um usuário pode estar associado a várias aplicações e relatórios, um aprendente pode possuir vários registros históricos e aplicações, e uma avaliação pode possuir diversas questões e ser utilizada em diversas aplicações. Cada aplicação, por sua vez, pode possuir várias respostas e, opcionalmente, um relatório. As questões também podem estar relacionadas a várias respostas, considerando diferentes aplicações do mesmo instrumento.

Para favorecer a consulta e o gerenciamento dos dados, foram definidos índices em campos utilizados com maior frequência nas operações do sistema. Entre eles estão o índice único para o e-mail dos usuários, índices relacionados ao status, nome e exclusão lógica dos aprendentes, status e categoria das avaliações, além de índices para os identificadores utilizados nas aplicações e respostas.

A Figura \ref{fig:diagrama_entidade_relacionamento_do_sistema} apresenta o Diagrama Entidade-Relacionamento desenvolvido para o sistema, permitindo visualizar graficamente as entidades, seus atributos e os relacionamentos estabelecidos entre elas.

\begin{figure}[h]
    \centering
    \includegraphics[width=0.9\textwidth]{der_sistema.png}
    \caption{Diagrama Entidade-Relacionamento do sistema}
    \label{fig:diagrama_entidade_relacionamento_do_sistema}
\end{figure}
